% Part: computability
% Chapter: recursive-functions
% Section: halting-problem

\documentclass[../../../include/open-logic-section]{subfiles}

\begin{document}

\olfileid{cmp}{rec}{hlt}
\olsection{Masalah Mandheg}

\emph{Masalah mandheg} umume yaiku masalah nemtokake,
yen diwenehi katrangan~$e$ (umpamane, program) ngenani fungsi
komputabel lan wilangan~$n$, apa komputasi fungsi kasebut
ing input~$n$ mandheg, yaiku ngasilake sawenehing asil.
Alan Turing mbuktekake yen masalah iki dhewe ora bisa
dirampungake dening fungsi komputabel, yaiku fungsi
\[
h(e, n) =
\begin{cases}
1 & \text{yen komputasi $e$ mandheg ing input $n$}\\
0 & \text{yen ora,}
\end{cases}
\]
ora komputabel.

Ing konteks fungsi rekursif parsial, peran
katrangan program bisa dilakoni dening indeks~$e$
kang diwenehake ing teorema wangun normal Kleene.
Yen $f$ fungsi rekursif parsial, saben $e$~kang
nyukupi persamaan ing teorema wangun normal iku
indeks saka~$f$. Yen diwenehi wilangan~$e$,
teorema wangun normal nyatakake yen
\[
\cfind{e}(x) \simeq U(\mu s \; T(e, x, s))
\]
iku rekursif parsial, lan kanggo saben fungsi
rekursif parsial $f\colon \Nat
\to \Nat$, ana $e \in \Nat$ kang ndadekake
$\cfind{e}(x) \simeq f(x)$ kanggo saben~$x \in \Nat$.
Malah, kanggo saben fungsi~$f$ kaya mangkono,
indeks~$e$ iku ora mung siji, nanging tanpa wates cacahé.
\emph{Fungsi mandheg}~$h$ ditegesi minangka
\[
h(e, x) =
\begin{cases}
1 & \text{yen $\cfind{e}(x) \fdefined$}\\
0 & \text{yen ora.}
\end{cases}
\]
Gatekna yen $h(e, x) = 0$ yen $\cfind{e}(x) \fundefined$,
nanging uga yen~$e$ dudu indeks saka fungsi
rekursif parsial babar pisan.

\begin{thm}
\ollabel{thm:halting-problem}
Fungsi mandheg~$h$ dudu fungsi rekursif parsial.
\end{thm}

\begin{proof}
Yen $h$ rekursif parsial, kita bisa netepake
\[
d(y) =
\begin{cases}
1 & \text{yen $h(y, y) = 0$}\\
\umin{x}{x \neq x} & \text{yen ora.}
\end{cases}
\]
Amarga ora ana wilangan $x$ kang nyukupi $x \neq x$,
mula $\umin{x}{x \neq x}$ ora ana, saengga
$d(y) \fundefined$ yen lan mung yen $h(y,y) \neq 0$.
Saka definisi iki kita entuk
\begin{enumerate}
\item $d(y) \fdefined$ yen lan mung yen
  $\cfind{y}(y) \fundefined$ utawa $y$
  dudu indeks saka fungsi rekursif parsial.
\item $d(y) \fundefined$ yen lan mung yen
  $\cfind{y}(y) \fdefined$.
\end{enumerate}
Yen $h$ rekursif parsial, $d$ uga rekursif parsial.
Mula, miturut teorema wangun normal Kleene,
fungsi iki nduwé indeks~$e_d$. Gatekna nilai
$h(e_d, e_d)$. Ana rong kemungkinan, $0$
lan~$1$.
\begin{enumerate}
\item Yen $h(e_d, e_d) = 1$, mula
  $\cfind{e_d}(e_d) \fdefined$. Nanging
  $\cfind{e_d} \simeq d$, lan $d(e_d)$
  ditegesi yen lan mung yen $h(e_d, e_d) = 0$.
  Dadi $h(e_d, e_d) \neq 1$.
\item Yen $h(e_d, e_d) = 0$, mula bisa uga
  $e_d$ dudu indeks saka fungsi rekursif parsial,
  utawa pancen indeks kasebut lan
  $\cfind{e_d}(e_d) \fundefined$. Nanging maneh,
  $\cfind{e_d} \simeq d$, lan $d(e_d)$ ora
  ditegesi yen lan mung yen
  $\cfind{e_d}(e_d) \fdefined$.
\end{enumerate}
Asile, $e_d$ jebul ora bisa dadi indeks
saka fungsi rekursif parsial. Nanging yen $h$
rekursif parsial, $d$ uga rekursif parsial,
mula panetepan $e_d$ minangka indekse iku sah.
Mula kita kudu nyimpulake yen $h$ ora bisa
dadi fungsi rekursif parsial.
\end{proof}

\end{document}
